\label{ch:detector-implementation}

This chapter develops every formula \texttt{tpeanuts.detector.icecube}
implements: the weighted-Monte-Carlo reweighting formula
(\cref{sec:reweighting}), the real 3-flavour Earth-matter oscillation
probability and its geometry (\cref{sec:earth-oscillation}), the real
atmospheric flux and detector-systematics hypersurface correction
(\cref{sec:flux-hypersurface}), and the free normalization parameters
(\cref{sec:free-normalization-icecube}).

\section{Weighted Monte Carlo reweighting}
\label{sec:reweighting}

\modulehead{detector/icecube/event\_rate.py}{\pyfunc{predicted\_neutrino\_counts}.}

\begin{theorybox}
For a real simulated event $i$ tagged (by its own PDG code, at the
interaction point) as final-state flavour $\beta_i$, the physically
correct expectation sums over every possible parent flavour $\alpha$ at
production, weighted by the real unoscillated atmospheric flux
$\Phi_\alpha$ and the oscillation transition probability
$P(\nu_\alpha\to\nu_{\beta_i})$:
\begin{equation}
  \keyresult{N_i(\theta) = \nu_{\rm norm}\cdot w_i \sum_{\alpha\in\{e,\mu\}}
    \Phi_\alpha(E_i,\cos\theta_{z,i})\,
    P(\nu_\alpha\to\nu_{\beta_i};\,E_i,\cos\theta_{z,i};\,\theta),}
  \label{eq:reweighting}
\end{equation}
with $\theta=(\theta_{23},\Delta m^2_{3l})$ the two free oscillation
parameters and $\nu_{\rm norm}$ a free overall normalization
(\cref{sec:free-normalization-icecube}). The sum runs over
$\alpha\in\{e,\mu\}$ only: the primary/conventional atmospheric
$\nu_\tau$ flux is negligible at these energies (the real flux table used
here, \cref{sec:flux-hypersurface}, has no $\nu_\tau$ entry either).
This applies identically to neutral-current and charged-current events --
an NC interaction still requires the neutrino to \emph{be} flavour
$\beta_i$ at that point. Every simulated event in a given real analysis
bin $k$ is summed to give that bin's predicted neutrino count,
\begin{equation}
  N_k^\nu(\theta) = \sum_{i\,\in\,{\rm bin}\,k} N_i(\theta).
  \label{eq:bin-sum}
\end{equation}
\end{theorybox}

\begin{implbox}
\pyfunc{predicted\_neutrino\_counts(theta23, deltam\_sq\_3l, ...,
channels, hypersurfaces)} loops over \code{CHANNELS = ("nc", "nue\_cc",
"numu\_cc", "nutau\_cc")}, calling
\pyfunc{tpeanuts.medium.earth.probability.earth\_probability\_transition}
once per channel (shape \code{(n\_events, 3, 3)}: [final $\beta$, initial
$\alpha$]) and reading off each event's own \emph{fixed} $\beta_i$ index
directly -- \labelcref{eq:reweighting}'s sum over $\alpha$, not $\beta$,
since $\beta_i$ is a per-event property, not summed. Binning
(\labelcref{eq:bin-sum}) is a differentiable \code{index\_add} into a
\code{(N\_BINS,)} tensor, verified (\texttt{test\_reweighting\_binning\_conserves\_total\_weight})
to conserve the total event weight exactly against an independently
computed direct per-event sum.
\end{implbox}

\section{Real 3-flavour Earth-matter oscillation}
\label{sec:earth-oscillation}

\modulehead{medium/earth/probability.py}{\pyfunc{earth\_probability\_transition}, shared unchanged with every other Earth-crossing detector in this project.}

\begin{theorybox}
A flavour state obeys $i\,d\psi/dx=H_f(x)\psi$, with the Standard-Model
flavour Hamiltonian
\begin{equation}
  H_f(x) = \frac{1}{2E}U\,{\rm diag}(0,\Delta m^2_{21},\Delta m^2_{31})\,U^\dagger
    + V_{\rm CC}(x)\,P_e, \qquad P_e={\rm diag}(1,0,0),
  \label{eq:flavour-hamiltonian}
\end{equation}
$U$ the PMNS matrix and $V_{\rm CC}(x)=\sqrt2\,G_F\,n_e(x)$ the Wolfenstein
charged-current matter potential \parencite{Wolfenstein1978}, sourced by
the Earth's electron density along the path. The exact evolution operator
is the path-ordered exponential
$S(x_2,x_1)=\mathcal T\exp[-i\int_{x_1}^{x_2}H_f\,dx]$; \texttt{tpeanuts}
evaluates it with a first-order perturbative expansion around each
traversed layer's mean density \parencite{Akhmedov2007Oscillograms}, the
same construction used for every Earth-matter effect in this project.

\textbf{Geometry.} Real analysis bins cover $\cos\theta_z\in[-1,\,0.1]$;
this is mapped to a detector nadir angle $\eta\in[0,\pi]$ via
\begin{equation}
  \eta = \arccos(-\cos\theta_z),
\end{equation}
so $\cos\theta_z=-1\Rightarrow\eta=0$ (straight through the Earth's centre)
and $\cos\theta_z=+1\Rightarrow\eta=\pi$ (no Earth crossing). With
$S\equiv S_{\rm Earth}(\eta,E;\theta)$,
\begin{equation}
  \keyresult{P(\nu_\alpha\to\nu_\beta;\,E,\eta;\,\theta) = \bigl|S_{\beta\alpha}(\eta,E;\theta)\bigr|^2.}
  \label{eq:transition-probability}
\end{equation}
Only $\theta_{23}$/$\Delta m^2_{3l}$ are free; $\theta_{12}$, $\theta_{13}$,
$\Delta m^2_{21}$, $\delta_{CP}$ are fixed at NuFIT~6.1
\parencite{EstebanNuFit2024} -- the standard treatment in every published
atmospheric-neutrino oscillation analysis, since $\theta_{23}$/$\Delta
m^2_{3l}$ overwhelmingly dominate the atmospheric $\nu_\mu$
disappearance/$\nu_\tau$ appearance signal.
\end{theorybox}

\begin{implbox}
\pyclass{AtmosphericOscillationModel} mirrors
\pyclass{VacuumOscillationModel}/\pyclass{SolarSMOscillationModel} but
frees $\theta_{23}$/$\Delta m^2_{3l}$ instead; its own
\code{.oscillation(theta, antinu=events.antinu)} takes a per-event boolean
antineutrino mask, since a real event sample mixes $\nu$ and $\bar\nu$
(unlike solar/reactor models, which fix a single sign).
\pyfunc{earth\_probability\_transition} implements
\labelcref{eq:transition-probability} unchanged from every other
Earth-crossing calculation in this project.
\end{implbox}

\section{Real atmospheric flux and detector-systematics hypersurface}
\label{sec:flux-hypersurface}

\modulehead{source/atmosphere/, detector/icecube/event\_rate.py}{\pyfunc{load\_atmospheric\_flux}, \pyfunc{interpolate\_hypersurface}.}

\begin{theorybox}
$\Phi_\alpha(E,\cos\theta_z)$ is the real Honda atmospheric flux
calculation \parencite{Honda2015Flux}, azimuth-averaged (this release's
own public binning has no azimuth axis) and bilinearly interpolated in
$(\log_{10}E,\cos\theta_z)$ at each event's own true kinematics -- a fixed
input, independent of $\theta$.

The release additionally provides real per-bin ``hypersurface''
corrections for 5 detector-calibration parameters $p_n$ (DOM efficiency,
hole-ice $p_0$/$p_1$, bulk-ice absorption/scattering), tabulated at 20 real
$\Delta m^2_{31}$ slices (event migration between reconstructed bins
itself depends weakly on the assumed $\Delta m^2_{31}$). At a fixed
systematics point $\{p_n\}$,
\begin{equation}
  \keyresult{C_k(p_1,\dots,p_5) = c_k + \sum_{n=1}^5 m_{k,n}\,(p_n-p_n^{\rm nominal}),}
  \label{eq:hypersurface}
\end{equation}
with $c_k$/$m_{k,n}$ the release's own real tabulated intercept/slopes,
piecewise-linearly interpolated (differentiably, in $\Delta m^2_{3l}$)
between the two tabulated slices bracketing the fitted value. This
document's forward model holds $\{p_n\}$ fixed at IceCube's own real
published best-fit point (\cref{tab:bestfit-systematics}) rather than
re-fitting them as free nuisances -- a deliberate scope decision, audited
directly in \cref{ch:results}.
\end{theorybox}

\begin{table}[htb]
\centering
\caption{Real detector-systematics reference points (Table IV of \cite{Abbasi2023IceCube}).}
\label{tab:bestfit-systematics}
\begin{tabular}{@{}lrr@{}}
\toprule
Parameter & Nominal & Best fit (used throughout) \\
\midrule
DOM efficiency & 1.00 & 1.06 \\
Hole-ice $p_0$ & 0.10 & $-0.27$ \\
Hole-ice $p_1$ & $-0.05$ & $-0.04$ \\
Bulk-ice absorption & 1.00 & 0.97 \\
Bulk-ice scattering & 1.00 & 0.99 \\
\bottomrule
\end{tabular}
\end{table}

\begin{implbox}
\pyfunc{interpolate\_hypersurface(table, deltam3l)} implements
\labelcref{eq:hypersurface}'s slice interpolation: exact at tabulated
points, linear at interior points, clamped outside the tabulated range
(verified directly, \texttt{test\_hypersurface\_interpolation\_matches\_tabulated\_slices\_and\_interior\_point}).
The full per-bin prediction composes
\labelcref{eq:reweighting,eq:bin-sum,eq:hypersurface} with the free muon
background normalization (\cref{sec:free-normalization-icecube}):
\begin{equation}
  N_k(\theta,\nu_{\rm norm},\mu_{\rm norm}) = \nu_{\rm norm}\cdot N_k^\nu(\theta)\cdot C_k(\Delta m^2_{3l})
    + \mu_{\rm norm}\cdot N_k^{\mu,{\rm MC}}.
  \label{eq:full-prediction}
\end{equation}
\end{implbox}

\section{Free normalization parameters}
\label{sec:free-normalization-icecube}

\begin{theorybox}
Unlike every other detector model in this project, $\theta$ here is
\emph{not} oscillation parameters alone: the release's own \code{weight}
column has no independently usable absolute livetime, so both a neutrino
($\nu_{\rm norm}$) and a muon-background ($\mu_{\rm norm}$) normalization
scale are always free, genuine fit parameters --
\begin{equation}
  \keyresult{\theta_{\rm fit} = (\theta_{23},\,\Delta m^2_{3l},\,\nu_{\rm norm},\,\mu_{\rm norm}),}
\end{equation}
4 free parameters always, regardless of any optional-flag convention used
elsewhere in this project. At $\nu_{\rm norm}=1$, the neutrino signal is
only $\sim$0.2\% of the total predicted rate (the muon background alone
dominates) -- $\nu_{\rm norm}=1$ is a wildly unnatural starting point for a
gradient-based fit, addressed directly in \cref{ch:results}'s
normalization-calibration step.
\end{theorybox}

\begin{implbox}
\pyclass{IceCubeDetectorModel.free} $=$ \code{("theta23", "DeltamSq3l",
"nu\_norm", "mu\_norm")} unconditionally. \code{.predict(theta)} unpacks
all four and implements \labelcref{eq:full-prediction} directly, verified
(\texttt{test\_gradient\_flows\_from\_all\_four\_free\_parameters}) to
carry a finite, non-zero autograd gradient with respect to every one of
them.
\end{implbox}
